% 評価

本章では、AnimatorSBTを3つの観点から評価する：
コストシミュレーション（cost simulation）による経済的実現可能性（\cref{subsec:cost_sim}）、
スケーラビリティ分析（\cref{subsec:scalability}）、
および\cref{subsec:formalization}で定義した要件に対する定性的評価である。

% ----- 6.1 コストシミュレーション -----
\subsection{コストシミュレーション}
\label{subsec:cost_sim}

AnimatorSBTの年間運用コストを、単一スタジオによるパイロット導入から業界全体での完全導入まで、5つの導入シナリオにわたってシミュレーションした。
本シミュレーションでは、労働力パラメータにJAniCA~(2023)およびNAFCA~(2024)の業界統計を使用し、ガスコストはPolygon PoS上のHardhatテストスイートから\textbf{実測}した値（ガス価格30~gwei、POL = \$0.22、2026年3月時点のレート）を用いている。

\subsubsection{パラメータ}

本シミュレーションでは、公開データに基づき以下の業界パラメータを仮定する：

\begin{itemize}
  \item アニメーション制作者総数：約5{,}500人（JAniCA 2023年調査ベース）
  \item フリーランス比率：47.3\%--69.6\%（JAniCA 2023 / 2019）
  \item アニメーター1人あたりの年間参加プロジェクト数：6
  \item バッチサイズ10でのバッチミント（デフォルト）
  \item 失効率（revocation rate）：発行済みSBTの1\%
  \item IPFSピンニング：500ファイル以下は無料枠、それ以上は\$20/月~\cite{pinata2026pricing}
\end{itemize}

\subsubsection{結果}

\Cref{tab:cost_scenarios}にシミュレーション結果を示す。

\begin{table}[ht]
\centering
\caption{5つの導入シナリオにおける年間コストシミュレーション。
すべてのコストはPolygon PoSレート（30 gwei、POL = \$0.22、2026年3月）でのUSD表記。
ガス値はHardhatテストスイートから実測した値である。}
\label{tab:cost_scenarios}
\footnotesize
\setlength{\tabcolsep}{3pt}
\begin{tabular}{@{}lrrrr@{}}
\toprule
\textbf{シナリオ} & \textbf{人数} & \textbf{SBT/年} & \textbf{年間コスト} & \textbf{1人あたり} \\
\midrule
Pilot (1 studio)      & 50    & 300      & \$0.22  & \$0.004 \\
Early (10 studios)    & 500   & 3{,}000  & \$242   & \$0.484 \\
Moderate (50\%)       & 1{,}300 & 7{,}800 & \$245  & \$0.189 \\
Full (low est.)       & 2{,}601 & 15{,}606 & \$251 & \$0.096 \\
Full (high est.)      & 3{,}827 & 22{,}962 & \$256 & \$0.067 \\
\bottomrule
\end{tabular}
\end{table}

\Cref{fig:cost_simulation}にコスト内訳と規模の経済性（economies of scale）を可視化する。

\begin{figure}[t]
  \centering
  \includegraphics[width=\columnwidth]{figures/cost_simulation}
  \caption{コストシミュレーション結果。左：5つの導入シナリオにおけるコンポーネント別年間コスト内訳。右：導入規模の拡大に伴うアニメーター1人あたりコストの推移（規模の経済性）。}
  \label{fig:cost_simulation}
\end{figure}

結果から顕著な知見が得られた：
\textbf{業界全体での完全導入（日本のフリーランスアニメーター全員）においても、年間運用コストの合計は約\$256}であり、東京での夕食1回分にも満たない金額である。
これはアニメ制作産業の年間収益3{,}621億円（\$24.1億）~\cite{tdb2025anime}の$0.000011\%$に相当する。

Pilotシナリオ（\$0.22）からEarlyシナリオ（\$242）への急激なコスト増加は、ブロックチェーンコストではなくIPFSピンニング料金の閾値（しきいち）に起因する：
Pilotの300 SBTはPinataの無料枠（500ファイル以下）に収まるが、Earlyシナリオの3{,}000 SBTでは有料プラン（\$20/月、\$240/年）が必要となる~\cite{pinata2026pricing}。
より広い観点では、コストの大部分はIPFSピンニング料金が占めており、ブロックチェーンのトランザクションコストではない。
完全導入時でもオンチェーンのミントコストはわずか\$16であり、Layer~2ブロックチェーンがガス料金をトランザクション単位では経済的に無視できる水準まで引き下げたことが確認された。

アニメーター1人あたりのコストは、固定的なIPFSコストがより多くのユーザーに分散されるため、規模の拡大とともに減少する。
完全導入時のコストはアニメーター1人あたり年間約\$0.067であり、中央値の時給1{,}111円~\cite{nafca2024labor}における1分未満の労働に相当する。

% ----- 6.2 スケーラビリティ -----
\subsection{スケーラビリティ分析}
\label{subsec:scalability}

\subsubsection{トランザクションスループット}

Polygon PoSは1日あたり約300万--800万トランザクションを処理し、ブロック生成時間は約2秒である~\cite{gangwal2023layer2}。
完全導入時、AnimatorSBTが生成するミントトランザクションは年間約22{,}962件、すなわち1日あたり約63件である。
これはPolygonの1日あたりの処理能力の0.002\%未満であり、システムがネットワークの限界を大きく下回る範囲で動作することを示している。

\subsubsection{ストレージ}

各SBTメタデータJSONファイルは約500バイトである。
完全導入時の年間メタデータストレージは
$22{,}962 \times 500 = 11.5$~MBであり、IPFSにとって無視できるほど小さい。
オンチェーンストレージは最小限（トークンID、オーナーアドレス、URI文字列）であり、トークンあたり約200バイトである。

\subsubsection{オンチェーンクエリコスト}

\texttt{tokensOfOwner()}関数は、発行済み全トークンに対して線形走査（linear scan）を行う（$O(n)$、$n$は総トークン数）。
現在の規模（数万トークン）では許容範囲であるが、$n$の増大に伴い非効率となる：
100{,}000トークンでは、\texttt{eth\_call}クエリに対するRPCノードのタイムアウト制限を超過する可能性がある。
本番環境での展開では、$O(1)$ルックアップのためにオーナーからトークンIDへのマッピングを追加するか、ポートフォリオクエリにはオフチェーンインデックス（例：The Graph）に依拠すべきである。

\subsubsection{成長予測}

アニメ市場の年平均成長率（CAGR）10.6\%~\cite{mordor2026anime}および比例的な労働力成長を仮定すると、
2031年までにシステムは年間約42{,}000 SBT（1日あたり約115件）を処理する必要がある。
これは依然としてPolygonの処理能力の範囲内であり、年間コストは約\$300--400と見込まれる。

% ----- 6.3 要件評価 -----
\subsection{要件充足性}
\label{subsec:requirements_eval}

\cref{subsec:formalization}で定義した5つの要件に対するAnimatorSBTの評価を以下に示す：

\begin{itemize}
  \item \textbf{R1（完全性）：} \textit{充足。}
    PMCを使用するすべてのアニメーターは、下請け階層における位置に関係なくSBTを受領できる。
    発行はPMC内の作業ログデータをトリガーとしており、
    スタジオの承認やクレジット委員会の判断には依存しない。

  \item \textbf{R2（検証可能性）：} \textit{充足。}
    第三者はパブリックブロックチェーンへのクエリによりSBTを検証できる。
    \texttt{evidence\_hash}は基礎となる作業ログデータへの暗号学的リンクを提供する。
    元のスタジオへの問い合わせは不要である。

  \item \textbf{R3（永続性）：} \textit{部分的に充足。}
    オンチェーンデータ（トークンの所有権、発行タイムスタンプ）はPolygonネットワークが稼働する限り永続する。
    しかし、IPFS上に保存された詳細メタデータは、少なくとも1つの当事者による継続的なピンニングに依存する。
    ピンニングサービスのサブスクリプションが失効し、他のノードがコピーを保持していない場合、
    オンチェーンレコードは維持されるもののメタデータにはアクセスできなくなる。
    これはIPFSベースのストレージの既知の制約であり、
    緩和策としてマルチプロバイダーピンニング、コミュニティ運営のIPFSノード、
    およびフォールバックとしてのArweaveベースの永続ストレージが挙げられる。

  \item \textbf{R4（ポータビリティ）：} \textit{充足。}
    SBTはアニメーターのウォレットアドレスに紐付けられ、
    アカウント抽象化（account abstraction）を通じてアニメーター自身が管理する。
    アニメーターは特定のプラットフォームに依存することなく、
    任意の検証者にウォレットアドレス（または派生識別子）を提示できる。

  \item \textbf{R5（最小限の摩擦）：} \textit{充足。}
    アニメーターのワークフローに変更はない---従来通りPMCに作業を記録し続けるだけである。
    SBTの発行はプロジェクトのマイルストーン到達時に自動的に行われる。
    ウォレット管理、ガス料金の支払い、ブロックチェーンに関する知識は不要である。
\end{itemize}

% ----- 6.4 制約事項 -----
\subsection{制約事項}
\label{subsec:eval_limitations}

現行の評価にはいくつかの制約事項があり、今後の研究で対処すべきである：

\begin{itemize}
  \item \textbf{テストネットのみの展開：}
    プロトタイプはPolygon Amoyテストネット上に展開されており、
    メインネットには展開されていない。
    ガスコストは正確に計測されている（EVM実行は同一）が、
    本番利用にはメインネットへの展開が必要である。
  \item \textbf{ユーザースタディの未実施：}
    現役のアニメーターを対象とした正式なユーザースタディはまだ実施していない。
    コストシミュレーションと実測ガスデータにより定量的な検証は行われているが、
    実環境でのユーザビリティと認知的価値は今後評価する必要がある。
  \item \textbf{ガス価格の変動性：}
    コスト推定はガス価格30~gweiおよびPOL価格\$0.22を固定値として仮定している。
    Polygon PoSのガス価格は歴史的に安定しているが、
    大幅な価格変動はコスト予測に影響を及ぼし得る。
  \item \textbf{自己申告データ：}
    SBTメタデータはPMC内の自己申告による作業ログから導出される。
    共同署名メカニズム（例：スタジオによる確認）がなければ、
    システムは貢献主張（contribution claims）の真実性を独立して検証することはできず、
    その存在と完全性のみを保証する。
    信頼モデルと緩和策については\cref{sec:discussion}で議論する。
  \item \textbf{IPFSメタデータの永続性：}
    SBTの詳細メタデータはIPFS上に存在し、アクセス可能な状態を維持するには
    アクティブなピンニングが必要である。すべてのピンニングプロバイダーが
    データのホスティングを停止した場合、オンチェーンのトークンレコードは残存するが、
    貢献の詳細は復元不可能となる。
  \item \textbf{導入に関する仮定：}
    シミュレーションでは業界全体での均一な導入を仮定している。
    実際にはネットワーク効果、スタジオの方針、文化的要因により、
    導入は不均一となる。
\end{itemize}
